% !TeX root = part8.tex
\documentclass[../final.tex]{subfiles}

\begin{document}

% Продолжение по плану семинаров: тема 6; сборник, задачи 7.2, 7.1, 7.4.
% 31.10.2026 занятия не было; отдельная часть за эту дату не создаётся.
\practice{8}{24.10.2026}{Электрон в периодическом поле}

\rtheory{Теорема Блоха и два предельных приближения}

Для потенциала $U(x+a)=U(x)$ собственные функции можно выбрать в виде
\[
    \psi_q(x)=e^{iqx}u_q(x),\qquad u_q(x+a)=u_q(x).
\]
Здесь $q$ --- \rconcept{квазиволновое число}, $\hbar q$ ---
квазиимпульс. Значения $q$ и $q+2\pi n/a$ эквивалентны;
выбираем первую зону $-\pi/a\le q<\pi/a$.
Внутри одной зоны спектр состоит из ветвей $E_\nu(q)$.

В \rconcept{приближении почти свободных электронов} периодический
потенциал мал по сравнению с характерной кинетической энергией.
Вдали от вырождения он даёт малые поправки к плоским волнам,
а около брэгговских плоскостей смешивает вырожденные состояния.

В \rconcept{приближении сильной связи} сначала решают задачу
об отдельной яме. Слабый переход между ямами превращает каждый
локальный уровень в узкую энергетическую зону.

\begingroup
\renewcommand{\thesection}{7.\arabic{section}}
\task[2]{Электрон в косинусоидальном потенциале}
\label{task:fksv_mathieu}

\condition
Для электрона массы $m$ в потенциале
\[
    U(x)=U_0\cos gx,\qquad g=\frac{2\pi}{a},
\]
найти энергетический спектр и волновые функции в приближениях
почти свободных электронов и сильной связи.

\solution

\subsection*{Уравнение для фурье-компонент}

Разложим функцию Блоха по плоским волнам:
\[
    \psi_q(x)=\frac1{\sqrt L}\sum_{\ell\in\mathbb Z}
       C_\ell e^{i(q+\ell g)x},\qquad
    \varepsilon_k=\frac{\hbar^2k^2}{2m}.
\]
У косинуса есть только две ненулевые фурье-компоненты
$U_{\pm g}=U_0/2$, поэтому уравнение Шредингера даёт
\[
    (\varepsilon_{q+\ell g}-E)C_\ell
       +\frac{U_0}{2}(C_{\ell-1}+C_{\ell+1})=0.
\]
Эта бесконечная система определяет зоны при любом $U_0$.
Для аналитического решения рассмотрим два предела.

\subsection*{Слабый потенциал вдали от границ зоны}

Введём энергетический масштаб $E_g=\hbar^2g^2/(2m)$ и считаем
$|U_0|\ll E_g$. Если состояние с волновым числом $k$
не вырождено с $k\pm g$, то первая поправка к энергии равна нулю:
$U_{0\text{ (Фурье)}}=0$. Вторая поправка равна
\begin{align*}
    E(k)&=\varepsilon_k+\frac{U_0^2}{4}
       \left(\frac1{\varepsilon_k-\varepsilon_{k-g}}
             +\frac1{\varepsilon_k-\varepsilon_{k+g}}\right)+\ldots\\
       &=\varepsilon_k+
          \frac{mU_0^2}{\hbar^2(4k^2-g^2)}+\ldots.
\end{align*}
Соответствующая функция с точностью до нормировки имеет вид
\[
    \psi_k(x)\simeq\frac1{\sqrt L}\left[
       e^{ikx}+\frac{U_0/2}{\varepsilon_k-\varepsilon_{k-g}}e^{i(k-g)x}
       +\frac{U_0/2}{\varepsilon_k-\varepsilon_{k+g}}e^{i(k+g)x}
    \right].
\]
Критерий применимости --- малость каждого отношения
$|U_0|/(2|\varepsilon_k-\varepsilon_{k\pm g}|)$.
Полюса формулы при $k=\pm g/2$ означают необходимость
вырожденной теории возмущений, а не бесконечную энергию.

\subsection*{Открытие щели на границе зоны}

Около $q=g/2$ положим $q=g/2+\kappa$, $|\kappa|\ll g$.
Сильно смешиваются две волны с числами $q$ и $q-g$:
\[
    H_{\mathrm{eff}}=
    \begin{pmatrix}
        \varepsilon_q&U_0/2\\
        U_0/2&\varepsilon_{q-g}
    \end{pmatrix}.
\]
Собственные значения этой матрицы равны
\begin{align*}
    E_\pm(q)&=\frac{\varepsilon_q+\varepsilon_{q-g}}2
       \pm\sqrt{\left(\frac{\varepsilon_q-\varepsilon_{q-g}}2\right)^2
                     +\frac{U_0^2}{4}}\\
       &=\frac{\hbar^2g^2}{8m}+\frac{\hbar^2\kappa^2}{2m}
          \pm\sqrt{\left(\frac{\hbar^2g\kappa}{2m}\right)^2
                         +\frac{U_0^2}{4}}.
\end{align*}
На самой границе
\[
    E_\pm(g/2)=\frac{\hbar^2g^2}{8m}\pm\frac{|U_0|}{2},
    \qquad \Delta=|U_0|.
\]
При $U_0>0$ нижнее состояние пропорционально
$\sin(gx/2)$, верхнее --- $\cos(gx/2)$:
\[
    \psi_-(x)=\sqrt{\frac2L}\sin\frac{gx}{2},\qquad
    \psi_+(x)=\sqrt{\frac2L}\cos\frac{gx}{2}.
\]
Не существенный общий фазовый множитель опущен.
Плотность нижнего состояния максимальна около минимумов потенциала.
При $U_0<0$ роли синуса и косинуса меняются.

Вдали от самой границы коэффициенты двух волн можно найти из
\[
    \frac{C_{-1}}{C_0}
       =\frac{E_\pm-\varepsilon_q}{U_0/2},\qquad
    |C_0|^2+|C_{-1}|^2=1.
\]
Двухволновое приближение описывает первую щель с точностью
до поправок от удалённых плоских волн.

\subsection*{Глубокие ямы: локальные уровни}

В пределе $|U_0|\gg E_g$ выбираем минимум $x_j=x_0+ja$,
где $U(x_j)=-|U_0|$. Вблизи него
\[
    U(x_j+\xi)=-|U_0|+\frac{|U_0|g^2}{2}\xi^2+O(\xi^4).
\]
Частота внутриминумного осциллятора и его длина равны
\[
    \Omega=g\sqrt{\frac{|U_0|}{m}},\qquad
    \ell_0=\sqrt{\frac{\hbar}{m\Omega}}\ll a.
\]
Низкие локальные уровни в гармоническом приближении:
\[
    \epsilon_\nu\simeq-|U_0|+\hbar\Omega\left(\nu+\frac12\right).
\]
Функция основного состояния в одной яме приближённо гауссова:
\[
    \varphi_0(\xi)=\frac1{\pi^{1/4}\sqrt{\ell_0}}
       e^{-\xi^2/(2\ell_0^2)}.
\]
Для возбуждённых состояний добавляется полином Эрмита.
Гармоническое приближение пригодно для уровней, существенно
удалённых от вершины барьера; оно не задаёт точный хвост функции
под всем барьером.

\subsection*{Блоховские комбинации и ширина зоны}

Обозначим ортонормированные локализованные функции выбранной зоны
через $w_\nu(x-x_j)$. С точностью до малых межъямных поправок
\[
    \psi_{\nu q}(x)=\frac1{\sqrt N}\sum_j
       e^{iqja}w_\nu(x-x_j).
\]
Удерживая только матричный элемент между соседними ямами
$t_\nu=\langle w_{\nu,j}|H|w_{\nu,j+1}\rangle$, получаем
\[
    E_\nu(q)\simeq\epsilon_\nu+2t_\nu\cos qa,
    \qquad W_\nu=4|t_\nu|.
\]
Для нижней зоны при обычном выборе фаз $t_0<0$, и её минимум
находится при $q=0$. Величина $|t_\nu|$ экспоненциально мала:
\[
    |t_\nu|\propto
       \exp\left[-\frac1\hbar\int_{x_1}^{x_2}
          \sqrt{2m\,[U(x)-\epsilon_\nu]}\,dx\right],
\]
где $x_1,x_2$ --- точки поворота, ограничивающие барьер
между двумя соседними ямами. Формула задаёт ведущую экспоненту;
предэкспоненциальный множитель требует сшивки решений.

\answer
В слабом потенциале спектр почти параболический, а на первой
границе зоны открывается щель $\Delta=|U_0|$.
В глубоком потенциале каждый локальный уровень образует узкую зону
$E_\nu(q)\simeq\epsilon_\nu+2t_\nu\cos qa$;
волновые функции --- блоховские суммы локализованных орбиталей.

\Needspace{19\baselineskip}
\task[1]{Модель Кронига--Пенни: слабая и сильная связь}
\label{task:fksv_kronig_penney}

\condition
Спектр электрона в потенциале $U(x)=\beta\sum_n\delta(x-na)$
определяется уравнением
\[
    \cos qa=\cos ka+\frac{m\beta}{\hbar^2k}\sin ka,
    \qquad E=\frac{\hbar^2k^2}{2m}.
\]
Получить приближения почти свободных электронов и сильной связи.
Здесь $q$ --- квазиволновое число, $k$ --- волновое число между барьерами;
в отличие от предыдущей задачи $\beta$ имеет размерность энергии,
умноженной на длину.

\solution

\subsection*{Условия сшивки и разрешённые зоны}

Функция непрерывна на каждом $\delta$-барьере, а её производная
испытывает скачок
\[
    \psi'(na+0)-\psi'(na-0)=\frac{2m\beta}{\hbar^2}\psi(na).
\]
Матрица переноса через ячейку имеет собственные значения $e^{\pm iqa}$.
Половина её следа и есть правая часть дисперсионного уравнения.
Для действительного $q$ разрешены энергии с
\[
    \left|\cos ka+\frac{m\beta}{\hbar^2k}\sin ka\right|\le1.
\]
Если модуль больше единицы, $q$ комплексно и энергия лежит в щели.

\subsection*{Малые барьеры}

Введём $P=m\beta a/\hbar^2$ и сначала считаем $|P|\ll1$.
При $\beta=0$ получаем свободный спектр, свёрнутый в первую зону.
Однако средняя фурье-компонента решётки дельта-функций ненулевая:
\[
    U_G=\frac\beta a\qquad \text{для всех }G=\frac{2\pi n}{a}.
\]
Поэтому вдали от вырождений первая поправка равна $\beta/a$,
а не нулю:
\[
    E(q)\simeq\frac{\hbar^2q^2}{2m}+\frac\beta a.
\]
Около первой границы, $q=\pi/a+\kappa$, двухволновая матрица даёт
\[
    E_\pm=E_B+\frac\beta a+\frac{\hbar^2\kappa^2}{2m}
       \pm\sqrt{\left(\frac{\hbar^2\pi\kappa}{ma}\right)^2
                         +\left(\frac\beta a\right)^2},
    \qquad E_B=\frac{\pi^2\hbar^2}{2ma^2}.
\]
Отсюда $\Delta=2|\beta|/a$.
Для слабых положительных барьеров края первой щели
равны $E_B$ и $E_B+2\beta/a$ с точностью первого порядка.

Тот же результат извлекается из точного уравнения. При $q=\pi/a$
положим $ka=\pi+\delta$. Тогда
\[
    -1\simeq-1+\frac{\delta^2}{2}-\frac P\pi\delta,
    \qquad \delta=0\quad\text{или}\quad\delta\simeq\frac{2P}{\pi}.
\]
Поскольку $E-E_B\simeq\pi\hbar^2\delta/(ma^2)$,
получаем ту же ширину щели.

\subsection*{Сильные положительные барьеры}

При $P\gg1$ ячейки почти изолированы. В нулевом приближении
это отрезки длины $a$ с непроницаемыми стенками:
\[
    k_na=n\pi,\qquad E_n^{(0)}=\frac{\hbar^2n^2\pi^2}{2ma^2},
    \qquad n=1,2,\ldots.
\]
Для фиксированного $n$ и достаточно большого $P$ положим
$ka=n\pi+\delta$, $|\delta|\ll1$. Тогда
\[
    \cos qa\simeq(-1)^n\left(1+\frac P{n\pi}\delta\right),
    \qquad
    \delta\simeq\frac{n\pi}{P}\bigl[(-1)^n\cos qa-1\bigr].
\]
Разлагая энергию, получаем
\begin{align*}
    E_n(q)&\simeq E_n^{(0)}+\frac{\hbar^2n\pi}{ma^2}\delta\\
       &=E_n^{(0)}+\frac{2E_n^{(0)}}P
          \bigl[(-1)^n\cos qa-1\bigr].
\end{align*}
Это косинусоидальная зона с параметрами
\[
    \bar E_n=E_n^{(0)}\left(1-\frac2P\right),\qquad
    t_n=\frac{(-1)^n E_n^{(0)}}P,\qquad
    W_n=\frac{4E_n^{(0)}}P.
\]
Для основной зоны $t_1<0$, минимум расположен при $q=0$.
Разложение неравномерно по номеру зоны: необходимо $n\pi/P\ll1$.
В отличие от гладкого высокого барьера, непрозрачность
$\delta$-барьера даёт степенное, а не экспоненциальное уменьшение ширины.

\subsection*{Сильные притягивающие центры}

Для $\beta<0$ отдельная дельта-яма имеет уровень
\[
    E_{\mathrm{at}}=-\frac{\hbar^2\kappa_0^2}{2m},\qquad
    \kappa_0=\frac{m|\beta|}{\hbar^2}.
\]
Для отрицательных энергий положим $k=i\kappa$. Точное уравнение принимает вид
\[
    \cos qa=\cosh\kappa a-\frac{\kappa_0}{\kappa}\sinh\kappa a.
\]
При $\kappa_0a\gg1$ ищем $\kappa=\kappa_0+\delta\kappa$:
\[
    \cos qa\simeq\frac{e^{\kappa_0a}}{2\kappa_0}\delta\kappa,
    \qquad
    \delta\kappa\simeq2\kappa_0e^{-\kappa_0a}\cos qa.
\]
Следовательно, с точностью ведущего экспоненциального порядка
\[
    E(q)\simeq E_{\mathrm{at}}
       -\frac{2\hbar^2\kappa_0^2}{m}e^{-\kappa_0a}\cos qa.
\]
Здесь сильная связь имеет обычный смысл перекрытия хвостов
локализованных атомных состояний.

\answer
При слабом потенциале первая щель равна $2|\beta|/a$.
При сильных положительных барьерах низкие зоны сужаются как $1/P$;
при глубокой цепочке притягивающих дельта-ям ширина нижней зоны
экспоненциально мала по $\kappa_0a$. В обоих случаях ведущая
дисперсия сильной связи косинусоидальна.
\endgroup

\Needspace{16\baselineskip}
\begingroup
\renewcommand{\thesection}{\arabic{section}}
\task[3]{Диагонализация гамильтониана сильной связи}
\label{task:fksv_tight_binding}

\condition
Получить спектр одноэлектронного гамильтониана
\[
    \hat H=E_0\sum_n\hat c_n^\dagger\hat c_n
       +t\sum_n(\hat c_n^\dagger\hat c_{n+1}
                         +\hat c_{n+1}^\dagger\hat c_n),
\]
где $t$ вещественно, а цепочка из $N$ узлов периодическая.
Обсудить учёт прыжков между следующими за ближайшими соседями.
Вывод для ближайших соседей соответствует задаче 7.4 сборника.

\Needspace{6\baselineskip}
\solution

\subsection*{Преобразование Фурье}

Введём фермионные операторы
\[
    \hat c_n=\frac1{\sqrt N}\sum_q e^{iqna}\hat c_q,
    \qquad
    \hat c_q=\frac1{\sqrt N}\sum_n e^{-iqna}\hat c_n,
    \qquad q=\frac{2\pi j}{Na}.
\]
Ортогональность $\sum_n e^{i(q-q')na}=N\delta_{qq'}$
сохраняет антикоммутаторы и даёт
\begin{align*}
    \sum_n\hat c_n^\dagger\hat c_n
       &=\sum_q\hat c_q^\dagger\hat c_q,\\
    \sum_n\hat c_n^\dagger\hat c_{n+1}
       &=\frac1N\sum_{q,q'}e^{iq'a}
          \hat c_q^\dagger\hat c_{q'}\sum_ne^{i(q'-q)na}
         =\sum_qe^{iqa}\hat c_q^\dagger\hat c_q.
\end{align*}
Добавляя эрмитово сопряжённый член, находим
\[
    \hat H=\sum_q E(q)\hat c_q^\dagger\hat c_q,
    \qquad E(q)=E_0+2t\cos qa.
\]
Одноэлектронное состояние имеет вид
\[
    |q\rangle=\frac1{\sqrt N}\sum_n e^{iqna}\hat c_n^\dagger|0\rangle.
\]
Его можно получить и из разностного уравнения
$(E-E_0)A_n=t(A_{n+1}+A_{n-1})$ подстановкой $A_n=e^{iqna}$.

\subsection*{Ширина зоны и эффективная масса}

Ширина зоны равна $4|t|$. При $t<0$ минимум находится при $q=0$:
\[
    E(q)\simeq E_0+2t-ta^2q^2
       =E_{\min}+\frac{\hbar^2q^2}{2m^*},\qquad
    m^*=\frac{\hbar^2}{2|t|a^2}.
\]
Групповая скорость
\[
    v(q)=\frac1\hbar\frac{dE}{dq}
       =-\frac{2ta}{\hbar}\sin qa
\]
обращается в нуль в центре и на границе зоны.

\subsection*{Что меняется при добавлении вторых соседей}

В ортонормированном узельном базисе допустим гамильтониан
\[
    \hat H_2=t_2\sum_n
       (\hat c_n^\dagger\hat c_{n+2}+\hat c_{n+2}^\dagger\hat c_n),
\]
который даёт
\[
    E(q)=E_0+2t\cos qa+2t_2\cos2qa.
\]
\rnote{Сам по себе этот член не приводит к неверному результату.}
В плане семинара сформулирован вопрос о некорректном учёте таких
переходов; его следует отличать от запрета модели со вторыми соседями.

Тонкость возникает, когда узельные состояния берут как
неортогональные орбитали изолированных атомов. Если
\[
    S_1=\langle\varphi_n|\varphi_{n+1}\rangle\ne0,
\]
то вместо обычной задачи на собственные значения получается
$HA=ESA$. В приближении соседних матричных элементов
\[
    E(q)=\frac{h_0+2h_1\cos qa}{1+2S_1\cos qa}.
\]
Даже здесь разложение по малому перекрытию создаёт вторую гармонику:
\begin{align*}
    E(q)&=h_0+2(h_1-h_0S_1)\cos qa
        +4(h_0S_1^2-h_1S_1)\cos^2qa+\ldots\\
       &=h_0+2(h_0S_1^2-h_1S_1)
        +2(h_1-h_0S_1)\cos qa\\
       &\hspace{1cm}+2(h_0S_1^2-h_1S_1)\cos2qa+\ldots.
\end{align*}
Для точности того же порядка нужно учитывать также собственные
матричные элементы и перекрытия вторых соседей, если они
сопоставимы с удержанными поправками.
Поэтому нельзя просто добавить малый прямой переход через два узла
и одновременно отбросить все поправки от неортогональности того же порядка.
После согласованной ортогонализации эффективные $E_0$, $t$, $t_2$
уже включают эти поправки.

\answer
Для ближайших соседей $E(q)=E_0+2t\cos qa$.
В ортонормированном базисе вторые соседи корректно добавляют
$2t_2\cos2qa$. Ошибка возникает при несогласованном приближении
в неортогональном атомном базисе, а не из-за самого наличия $t_2$.
\endgroup

\clearpage
\begingroup
\renewcommand{\thesection}{\arabic{section}}
\task[4]{Низкотемпературная электронная теплоёмкость в 1D и 2D}

\condition
Определить температурную зависимость теплоёмкости вырожденного
невзаимодействующего электронного газа в одном и двух измерениях.
Для явных коэффициентов использовать параболический спектр
$\varepsilon_k=\hbar^2k^2/(2m)$, спиновое вырождение $g_s=2$
и фиксированное число электронов $N_e$.

\solution

\subsection*{Разложение Зоммерфельда при фиксированном числе частиц}

Пусть $D(E)$ --- полная плотность одноэлектронных состояний,
включая спин. Тогда
\[
    N_e=\int_0^\infty D(E)f(E)\,dE,\qquad
    U=\int_0^\infty ED(E)f(E)\,dE,
    \qquad f(E)=\frac1{e^{(E-\mu)/(k_{\mathrm B}T)}+1}.
\]
Если $D(E)$ гладка в окрестности $E_F$ и
$k_{\mathrm B}T\ll E_F$, разложение Зоммерфельда даёт
\[
    \int_0^\infty F(E)f(E)\,dE
       =\int_0^\mu F(E)\,dE
          +\frac{\pi^2}{6}(k_{\mathrm B}T)^2F'(\mu)+\ldots.
\]
Из постоянства $N_e$ следует
\[
    \delta\mu=\mu-E_F
       =-\frac{\pi^2}{6}(k_{\mathrm B}T)^2
          \frac{D'(E_F)}{D(E_F)}+\ldots.
\]
Подставляя это в энергию, получаем сокращение членов с $D'(E_F)$:
\begin{align*}
    U(T)-U(0)
       &=E_FD(E_F)\delta\mu
          +\frac{\pi^2}{6}(k_{\mathrm B}T)^2
             [D(E_F)+E_FD'(E_F)]+\ldots\\
       &=\frac{\pi^2}{6}D(E_F)(k_{\mathrm B}T)^2+\ldots.
\end{align*}
Следовательно,
\[
    C_V=\frac{\pi^2}{3}D(E_F)k_{\mathrm B}^2T+o(T).
\]
Линейность связана с тем, что тепловые возбуждения затрагивают
лишь слой ширины $\sim k_{\mathrm B}T$ около поверхности Ферми.

\subsection*{Одномерный газ}

В отрезке длины $L$ число состояний с $|k|\le k_E$
равно $g_sLk_E/\pi$. Поэтому
\[
    D_{1\mathrm D}(E)=\frac{g_sL}{\pi\hbar}
       \sqrt{\frac m{2E}},\qquad
    D_{1\mathrm D}(E_F)=\frac{N_e}{2E_F}.
\]
Отсюда
\[
    C_{1\mathrm D}=\frac{\pi^2}{6}N_e k_{\mathrm B}
       \frac{k_{\mathrm B}T}{E_F}.
\]
Поскольку $D'/D=-1/(2E)$, химический потенциал в этом случае
слегка растёт с температурой:
\[
    \mu(T)=E_F+\frac{\pi^2}{12}
       \frac{(k_{\mathrm B}T)^2}{E_F}+\ldots.
\]
Для $g_s=2$, $v_F=\hbar k_F/m$ ту же теплоёмкость на длину
можно записать как
\[
    \frac{C_{1\mathrm D}}L=\frac{2\pi}{3}
       \frac{k_{\mathrm B}^2T}{\hbar v_F}.
\]

\subsection*{Двумерный газ}

Для площади $A$ число состояний внутри круга равно
$g_sAk_E^2/(4\pi)$. Следовательно,
\[
    D_{2\mathrm D}(E)=\frac{g_sAm}{2\pi\hbar^2},\qquad
    N_e=D_{2\mathrm D}E_F.
\]
Плотность состояний постоянна, и
\[
    C_{2\mathrm D}=\frac{\pi^2}{3}N_e k_{\mathrm B}
       \frac{k_{\mathrm B}T}{E_F},\qquad
    \frac{C_{2\mathrm D}}A=\frac{\pi m}{3\hbar^2}k_{\mathrm B}^2T
       \quad(g_s=2).
\]
Число частиц можно проинтегрировать точно:
\[
    N_e=D_{2\mathrm D}k_{\mathrm B}T
       \ln(1+e^{\mu/(k_{\mathrm B}T)}),
    \qquad
    \mu(T)=k_{\mathrm B}T\ln\left(e^{E_F/(k_{\mathrm B}T)}-1\right).
\]
Поэтому $\mu-E_F$ экспоненциально мало при $T\to0$;
поправки порядка $T^2$ в двумерном параболическом газе нет.

\subsection*{Применимость к зоне периодического кристалла}

Для произвольной зоны формула через $D(E_F)$ сохраняется,
если уровень Ферми находится в гладкой части спектра и
далеко от края зоны или особенности Ван-Хова на масштабе $k_{\mathrm B}T$.
Если он попал в запрещённую зону, металлический линейный закон
не применим. Для конечной цепочки температура должна также
превышать характерное расстояние между уровнями около $E_F$.

\answer
\rresult{И в 1D, и в 2D электронная теплоёмкость металла линейна по $T$.}
Для параболического газа
\[
    C_{1\mathrm D}=\frac{\pi^2}{6}N_e k_{\mathrm B}\frac T{T_F},
    \qquad
    C_{2\mathrm D}=\frac{\pi^2}{3}N_e k_{\mathrm B}\frac T{T_F},
    \qquad T_F=\frac{E_F}{k_{\mathrm B}}.
\]
В отличие от фононов с линейной дисперсией, степень температуры
здесь не равна размерности системы: существенен узкий слой около $E_F$.
\endgroup

\end{document}
